\chapter{Requirements Traceability}
\label{ch:traceability}

\section{Purpose}

This chapter demonstrates that the design follows from the requirements rather
than from implementation convenience. Each requirement from the System
Requirements Specification is listed with the diagram or diagrams that realise
it, and those not yet realised are identified so that the gaps are visible.

\section{Traceability Matrix}

\begin{longtable}{L{0.07\textwidth}L{0.52\textwidth}L{0.33\textwidth}}
	\caption{Requirements to diagrams}\label{tab:traceability}\\
	\toprule
	\textbf{ID} & \textbf{Requirement} & \textbf{Realised in} \\
	\midrule
	\endfirsthead
	\toprule
	\textbf{ID} & \textbf{Requirement} & \textbf{Realised in} \\
	\midrule
	\endhead

	\multicolumn{3}{l}{\textit{3.1.1 State Diagrams to Activity Diagrams}} \\
	\addlinespace
	\req{SA-1} & Read a state machine from GENESYS, including all states,
	transitions, triggers, conditions and entry and exit behaviour.
	& Figure~\ref{fig:sa-activity} \\
	\addlinespace
	\req{SA-2} & Create an activity diagram in the repository representing that
	state machine's behaviour.
	& Figures~\ref{fig:sa-activity}, \ref{fig:class} \\
	\addlinespace
	\req{SA-3} & Offer an option to show all possible states and transitions in
	activity diagram outputs.
	& Figure~\ref{fig:sa-activity} \\
	\addlinespace

	\multicolumn{3}{l}{\textit{3.1.2 Activity Diagrams to Code}} \\
	\addlinespace
	\req{AC-1} & Read an activity diagram and generate C\# or Visual Basic code
	implementing its control flow.
	& Figures~\ref{fig:export-activity}, \ref{fig:usecase}, \ref{fig:class} \\
	\addlinespace
	\req{AC-2} & Each activity node becomes a class, with function prototypes
	attached to each generated function class.
	& Figures~\ref{fig:export-activity}, \ref{fig:class} \\
	\addlinespace
	\req{AC-3} & Forks, loops and merges become the corresponding constructs of
	the target language.
	& Figure~\ref{fig:export-activity} \\
	\addlinespace
	\req{AC-4} & Generated code shall be compilable.
	& Figure~\ref{fig:export-activity} \\
	\addlinespace

	\multicolumn{3}{l}{\textit{3.1.3 Code to State Diagram}} \\
	\addlinespace
	\req{CS-1} & Load a code repository and extract a control flow graph.
	& Figures~\ref{fig:import-activity}, \ref{fig:block}, \ref{fig:class} \\
	\addlinespace
	\req{CS-2} & The user defines initial variable conditions, and translation
	begins there and follows calls to a designated depth.
	& Figure~\ref{fig:import-activity} \\
	\addlinespace
	\req{CS-3} & Output covers every possible branch when initial conditions are
	properly defined.
	& Figure~\ref{fig:import-activity} \\
	\addlinespace
	\req{CS-4} & Loops and recursive calls are represented as a single item,
	as states that call themselves with differing transitions, rather than
	expanded per occurrence.
	& Figure~\ref{fig:import-activity} \\
	\addlinespace
	\req{CS-5} & Unused variables, methods and classes are reported to the user.
	& Figure~\ref{fig:import-activity} \\
	\addlinespace
	\req{CS-6} & The tool derives a name for each generated state from
	information in the code.
	& Figure~\ref{fig:import-activity} \\
	\addlinespace

	\multicolumn{3}{l}{\textit{3.1.4 Stretch Functional Requirements}} \\
	\addlinespace
	\req{FS-1} & All functionality accessible from the GENESYS interface via
	scripting, with the application running as a background service.
	& Figures~\ref{fig:block}, \ref{fig:state}, \ref{fig:sequence} \\
	\addlinespace
	\req{FS-2} & A status window after use, showing where the model and the code
	have changed and whether the operation succeeded.
	& Figure~\ref{fig:state} \\
	\addlinespace
	\req{FS-3} & Renames and file moves reconciled rather than treated as a
	deletion and a creation.
	& Figure~\ref{fig:import-activity}, stretch \\
	\addlinespace

	\multicolumn{3}{l}{\textit{3.2.1 Non-Functional, minimum viable product}} \\
	\addlinespace
	\req{NFR-1} & Intuitive for engineers.
	& Figures~\ref{fig:usecase}, \ref{fig:state} \\
	\addlinespace
	\req{NFR-2} & A failure in the tool shall not terminate GENESYS.
	& Figure~\ref{fig:block} \\
	\addlinespace
	\req{NFR-3} & A cancelled or failed operation shall leave no changes
	partially written.
	& Figures~\ref{fig:state}, \ref{fig:import-activity} \\
	\addlinespace
	\req{NFR-4} & The same inputs and configuration shall always produce the same
	output.
	& Not visible in any diagram \\
	\addlinespace
	\req{NFR-5} & Windows only, matching GENESYS.
	& Not a diagram concern \\
	\addlinespace

	\multicolumn{3}{l}{\textit{3.2.2 Non-Functional, stretch}} \\
	\addlinespace
	\req{NFS-1} & Translation speed optimisation.
	& Not a diagram concern \\
	\addlinespace
	\req{NFS-2} & Remain responsive on large repositories.
	& Not a diagram concern \\
	\addlinespace
	\req{NFS-3} & Complete operations within a reasonable time for the
	translation size.
	& Not a diagram concern \\
	\bottomrule
\end{longtable}

\section{How Each Diagram Realises Its Requirements}

\paragraph{Block diagram, Figure~\ref{fig:block}.}
The separation between the \emph{In GENESYS} group and the \emph{.NET Process}
group is what satisfies \req{NFR-2}: a fault in the analysis process cannot
terminate the host, because they are different processes. The same separation is
what makes \req{FS-1} reachable without redesign, since an interface placed
inside GENESYS drives an application that already runs separately. The Code
Repository and Code Extractor realise the loading half of \req{CS-1}.

\paragraph{Use case diagram, Figure~\ref{fig:usecase}.}
\emph{Export Model to Code} and \emph{Import Model from Code} correspond to
\req{AC-1} and \req{CS-1} respectively. A single Engineer actor, rather than
separate roles per direction, reflects \req{NFR-1}: the same person performs both
translations through one interface.

\paragraph{State to activity diagram, Figure~\ref{fig:sa-activity}.}
The parallel reads satisfy \req{SA-1}, covering states with their entry and exit
behaviours alongside transitions with their triggers and conditions. The loop
body realises \req{SA-2}, with \emph{Create Decision Node per Guard} being the
step that turns a transition condition into a branch. \emph{Mark State as
Reached} and \emph{Append Unreached States and Transitions} together realise
\req{SA-3}, the second being possible only because the first records what the
traversal covered.

\paragraph{Import activity diagram, Figure~\ref{fig:import-activity}.}
\emph{Select User Configurations} and \emph{Enqueue Called Methods within Depth
Limit} together realise \req{CS-2}. \emph{Build Control Flow Graph} realises the
extraction half of \req{CS-1}, and \emph{Walk Control Flow Graph} with
\emph{Reducible to Constructs?} realises \req{CS-3}. The bounded queue satisfies
\req{CS-4}, since a method already processed is not enqueued again, so recursion
produces one item rather than an expansion. The \emph{Review and Write} node
satisfies \req{NFR-3}, because no path reaches a write without passing through
it. \emph{Derive State Names from Code} realises \req{CS-6}, and \emph{Report
Unused Variables, Methods and Classes} realises \req{CS-5}, placed after the
queue drains because reachability determines what counts as unused.
\emph{Reconcile Renames and Moves} covers \req{FS-3} and is drawn dashed, since
it is a stretch requirement.

\paragraph{Export activity diagram, Figure~\ref{fig:export-activity}.}
\emph{Get Functions} through \emph{Build Method} realise \req{AC-1} and
\req{AC-2}, with \emph{Relate to Owning Class} being the step that places each
generated prototype on its function class. \emph{Build Input/Output} produces the
prototypes themselves. The forked branches of \emph{Build Method} correspond to
the construct mapping required by \req{AC-3}.

\paragraph{Sequence diagram, Figure~\ref{fig:sequence}.}
The three lifelines show the application as a participant distinct from GENESYS,
which is the interaction \req{FS-1} would route through a script rather than
through the user directly.

\paragraph{Class diagram, Figure~\ref{fig:class}.}
\texttt{Function} carrying \texttt{parameters[]} is what makes \req{AC-2}
implementable, since a generated class needs its prototypes' signatures.
\texttt{Absyn} as the pivot between \texttt{C\#} and \texttt{Genesys Project}
is what allows both \req{AC-1} and \req{CS-1} to be served by one chain rather
than two independent translators.

\paragraph{State diagram, Figure~\ref{fig:state}.}
\emph{Showing Report}, reached from success, denial and failure alike, realises
\req{FS-2}. \emph{Awaiting Review} preceding \emph{Writing} realises
\req{NFR-3}, since the only route to a write is through acceptance.

\section{Requirements Not Yet Realised}

\begin{longtable}{L{0.10\textwidth}L{0.82\textwidth}}
	\caption{Gaps against the requirements}\label{tab:omissions}\\
	\toprule
	\textbf{ID} & \textbf{Why it is absent, and what is needed} \\
	\midrule
	\endfirsthead
	\toprule
	\textbf{ID} & \textbf{Why it is absent, and what is needed} \\
	\midrule
	\endhead

	\req{NFR-4} & That the same inputs always produce the same output is a
	property of the implementation rather than of the structure, so no diagram
	shows it. \\
	\addlinespace

	\req{NFR-5}, \req{NFS-1}, \req{NFS-2}, \req{NFS-3} & Platform and performance
	requirements are properties of the delivered system rather than of its
	structure or behaviour, and are addressed in the requirements specification
	rather than here. \\
	\bottomrule
\end{longtable}